\expandafter\def\csname HistoricalNoteText4\endcsname{In einem Briefe an \textsc{Leibnitz}. Horrebowii Opera T. II p. 142.}
\expandafter\def\csname HistoricalNoteText5\endcsname{Dieser ist also dem Sinus der mittleren Horizontal-Parallaxe der Sonne gleich, wenn die mittlere
Entfernung der Erde von der Sonne als Einheit angenommen wird.}
\clearpage\SourceLeaf{99}%
\pdfbookmark[1]{Einige Bemerkungen zur Vereinfachung der Rechnung für die geocentrischen Örter der Planeten}{bemerkungen1804}
\begin{center}
{\small EINIGE BEMERKUNGEN ZUR VEREINFACHUNG DER RECHNUNG
FÜR DIE GEOCENTRISCHEN ÖRTER DER PLANETEN.\par}
\medskip\hrule\medskip
{\small Monatliche Correspondenz zur Beförderung der Erd- und Himmels-Kunde,\par
herausgegeben vom Freiherrn \textsc{von Zach}. Mai 1804.\par}
\medskip\hrule
\end{center}\medskip
Seit der Erfindung der Pendeluhren beziehen sich alle unsere Beobachtungen
der Fixsterne, Planeten und Cometen nicht auf ihre Lage gegen die Ecliptik,
sondern unmittelbar auf ihre Lage gegen den Äquator. In unsern neuesten
und besten Sternverzeichnissen und Sternkarten sind gleichfalls nicht Länge und
Breite, sondern Rectascension und Declination zum Grunde gelegt. Man hat
daher sehr häufig Veranlassung, für Planeten und Cometen ihre geocentrischen
Örter in Beziehung auf den Äquator aus ihren heliocentrischen Örtern in ihrer
Bahn zu berechnen; und man würde diese Veranlassung noch häufiger haben,
wenn man sich entschlösse, in den astronomischen Ephemeriden anstatt der wenig
nutzenden Längen und Breiten der Planeten durchgängig die in jeder practischen
Hinsicht viel brauchbarern geraden Aufsteigungen und Abweichungen anzusetzen.
Dies hat der vortreffliche \textsc{Römer} bereits vor hundert Jahren angerathen\HistoricalNote{4},
und besonders wird es ganz unentbehrlich für die beiden neuesten Planeten,
die so schwer zu beobachten, und nur vermittelst sehr detaillirter Himmelskarten
aus den sie umgebenden kleinen Fixsternen herauszufinden sind. Eben
so häufig würde die allgemeinere Befolgung eines andern Vorschlages zu jener
Rechnung Gelegenheit geben, nemlich bei Vergleichung des beobachteten Orts
\SourceLeaf{100}%
eines Planeten oder Cometen mit dem berechneten unmittelbar die beobachtete
gerade Aufsteigung und Abweichung zum Grunde zu legen, und nicht erst, wie
gewöhnlich geschieht, aus diesen eine sogenannte beobachtete Länge und Breite
abzuleiten. Die mit diesem Verfahren verbundenen Vortheile sind bereits von
einem competenten Richter im V. Bande der Mon. Corr. S. 594 erwähnt worden.

Aus diesem Gesichtspunkte hat man die geocentrische Länge und Breite
des Planeten nur als Mittelgrössen anzusehen, um seine Lage gegen den Äquator
zu finden. Es wird daher obigen Vorschlägen vielleicht zu einer Empfehlung
mehr dienen, dass man dieser Zwischenrechnung, ja selbst der Reduction des
heliocentrischen Orts in der Bahn auf den heliocentrischen Ort in Beziehung auf
die Ecliptik ganz überhoben sein, und durch sehr einfache und geschmeidige
Formeln, welche in gegenwärtigem Aufsatze entwickelt werden sollen, aus jenem
die geocentrische Rectascension und Declination unmittelbar ableiten kann. Zu
diesen Vortheilen kann man noch die grosse Leichtigkeit hinzufügen, womit sich
bei diesem Verfahren die \textit{Parallaxe} auch in dem Falle mit in Rechnung bringen
lässt, wenn der Planet sich ausser dem Meridiane des Beobachtungsorts befindet,
welches zwar seltener nöthig, dann aber auch bei andern Methoden ungleich beschwerlicher
ist.

Durch den Mittelpunkt der Sonne lege man drei auf einander senkrechte
Ebenen, die eine parallel mit dem Erd-Äquator, die zweite durch die Punkte
der Nachtgleichen, also die dritte durch die Punkte der Sonnenwenden. Es
heissen die senkrechten Abstände des Mittelpunkts der Erde von diesen drei
Ebenen respective $Z$, $Y$, $X$, und die Abstände eines Planeten von eben denselben
$z$, $y$, $x$. Diese Abstände sollen als positiv angenommen werden bei der ersten
Ebene auf der Seite, wo der Nordpol liegt, bei der zweiten auf der Seite
der Sommer-Sonnenwende, bei der dritten auf der Seite der Frühlings-Nachtgleiche.
Es werden demnach $z-Z$, $y-Y$, $x-X$ die auf ähnliche Art genommenen
senkrechten Abstände des Planeten von dreien, den obigen parallel durch
den Mittelpunkt der Erde gelegten Ebenen sein. Bezeichnet man also die geocentrische
gerade Aufsteigung des Planeten durch $\alpha$, seine Abweichung durch
$\delta$, den Abstand von der Erde durch $\Delta$, so wird
\[
\begin{aligned}
x-X&=\Delta\cos\delta\cos\alpha\\
y-Y&=\Delta\cos\delta\sin\alpha\\
z-Z&=\Delta\sin\delta
\end{aligned}
\]
\SourceLeaf{101}%
Man findet folglich $\alpha$ durch die Formel $\operatorname{tang}\alpha=\frac{y-Y}{x-X}$, wo das positive oder
negative Zeichen des Zählers entscheiden muss, ob $\alpha$ in den beiden ersten oder
in den beiden letzten Quadranten anzunehmen ist. Sodann wird
\[
\Delta\cos\delta=\frac{x-X}{\cos\alpha}=\frac{y-Y}{\sin\alpha},\quad
\text{und}\quad\operatorname{tang}\delta=\frac{z-Z}{\Delta\cos\delta}
\]

Auf diese Weise erhält man also die Rectascension und Declination des
Planeten aus dem Mittelpunkte der Erde gesehen. Verlangt man dieselben, wie
sie aus einem Punkte auf der Oberfläche der Erde erscheinen, so ist in obigen
Formeln weiter keine Änderung nöthig, als dass man statt der Coordinaten des
Mittelpunkts $X$, $Y$, $Z$, die Abstände des Beobachtungsortes von den drei Fundamental-Ebenen
gebrauchen muss. Ist der Halbmesser der Erde $=\rho$\HistoricalNote{5}, die
Polhöhe des Beobachtungsorts $=\varphi$, und die Sternzeit, die derselbe im Augenblicke
der Beobachtung zählt, im Bogen, oder die gerade Aufsteigung des culminirenden
Punkts des Äquators $=\vartheta$: so werden jene Abstände, wie man leicht
übersehen wird:
\[
\begin{gathered}
X+\rho\cos\varphi\cos\vartheta\\
Y+\rho\cos\varphi\sin\vartheta\\
Z+\rho\sin\varphi
\end{gathered}
\]
Hiebei ist die Erde als eine Kugel angenommen. Fände man es nöthig, auch
auf die sphäroidische Gestalt der Erde Rücksicht zu nehmen (welcher Fall bei
Cometen eintreten könnte, die der Erde sehr nahe kämen), so dürfte man nur
für $\rho$ die Entfernung des Beobachtungsorts vom Mittelpunkte der Erde, und für
$\varphi$ seine sogenannte verbesserte Polhöhe setzen, die nach bekannten Regeln bestimmt
werden.

Man sieht jetzt also, dass es lediglich darauf ankommt, eine bequeme Methode
zur Bestimmung der Coordinaten $X$, $Y$, $Z$, $x$, $y$, $z$ aufzusuchen. In dieser
Absicht sei um die Sonne eine Kugelfläche mit unbestimmtem Halbmesser
beschrieben; auf derselben bezeichne $P$ den Nordpol der Ecliptik, $p$ den Nordpol
der Ebene der Planeten-Bahn; $K$ den Ort der Erde, $k$ den heliocentrischen
Ort des Planeten; endlich $\mathfrak{x}$, $\mathfrak{y}$, $\mathfrak{z}$ diejenigen Pole der drei Fundamental-Ebenen,
die auf der Seite liegen, wo die Abstände $x$, $y$, $z$ positiv genommen wer\SourceLeaf{102}%
den: also $\mathfrak{z}$ den Nordpol des Äquators, $\mathfrak{x}$ den Punkt der Frühlings-Nachtgleiche,
$\mathfrak{y}$ den Punkt des Äquators, der $90^0$ Rectascension hat (eine Figur wird
sich hienach jeder, der es nöthig findet, leicht selbst entwerfen können). Setzen
wir nun den Abstand der Erde von der Sonne $=R$, so wird offenbar
\[
\begin{aligned}
X&=R\cos\mathfrak{x}K\\
Y&=R\cos\mathfrak{y}K\\
Z&=R\cos\mathfrak{z}K
\end{aligned}
\]
Folglich, da in dem sphärischen Dreiecke $\mathfrak{x}PK$ die Seite $PK=90^0$, also
$\cos\mathfrak{x}K=\sin\mathfrak{x}P\cos\mathfrak{x}PK$ ist,
\[
\begin{aligned}
X&=R\sin\mathfrak{x}P\cos\mathfrak{x}PK,\quad\text{und eben so}\\
Y&=R\sin\mathfrak{y}P\cos\mathfrak{y}PK\quad\text{und}\\
Z&=R\sin\mathfrak{z}P\cos\mathfrak{z}PK
\end{aligned}
\]

Ganz auf ähnliche Weise werden die Coordinaten des Planeten, wenn wir
dessen Abstand von der Sonne durch $r$ bezeichnen,
\[
\begin{aligned}
x&=r\sin\mathfrak{x}p\cos\mathfrak{x}pk\\
y&=r\sin\mathfrak{y}p\cos\mathfrak{y}pk\\
z&=r\sin\mathfrak{z}p\cos\mathfrak{z}pk
\end{aligned}
\]

Wir bemerken hier ein für allemal, dass wir den sphärischen Winkel
$\mathfrak{x}PK$ so verstanden wissen wollen, wie der Schenkel $PK$ auf den Schenkel
$P\mathfrak{x}$ \textit{nach der Ordnung der Zeichen folgt}, so dass also derselbe mit $KP\mathfrak{x}$ nicht
gleichbedeutend sein soll, sondern beide einander zu $360^0$ ergänzen. Eben so
soll jeder andere sphärische Winkel zu verstehen sein. Durch eine solche nähere
Bestimmung gewinnen wir den Vortheil, dass die Grundformeln der sphärischen
Trigonometrie sich ohne weiteres auch auf Dreiecke mit Winkeln über
$180^0$ ausdehnen lassen, und weichen so der sonst Statt findenden Nothwendigkeit
aus, mehrere einzelne Fälle unterscheiden zu müssen. Übrigens werden
Winkel, deren Unterschied $360^0$ oder ein Vielfaches davon beträgt, jederzeit
als gleichbedeutend angesehen werden.

Wir nehmen nun zuvörderst die Coordinaten $X$, $Y$, $Z$ vor, und setzen die
Schiefe der Ecliptik $=\varepsilon$, die heliocentrische Länge der Erde $=\lambda$ ($=$ geocentrische
Länge der Sonne $+180^0$). In obigen Formeln wird also $\mathfrak{x}P=90^0$,
$\mathfrak{y}P=90^0+\varepsilon$, $\mathfrak{z}P=\varepsilon$, $\mathfrak{x}PK=\lambda$, $\mathfrak{y}PK=\mathfrak{z}PK=\lambda-90^0$, folglich
\SourceLeaf{103}%
\[
\begin{aligned}
X&=R\cos\lambda\\
Y&=R\sin\lambda\cos\varepsilon\\
Z&=R\sin\lambda\sin\varepsilon
\end{aligned}
\]

Für den Planeten setzen wir Kürze halber $\mathfrak{x}p=a$, $\mathfrak{y}p=b$, $\mathfrak{z}p=c$, seine
Entfernung in der Bahn vom aufsteigenden Knoten auf der Ecliptik $=t$, und
die Winkel $\mathfrak{x}pP$, $\mathfrak{y}pP$, $\mathfrak{z}pP$ respective $=A$, $B$, $C$. Man wird leicht übersehen,
dass $Ppk=t-90^0$ (oder nach obiger Anmerkung $=t+270^0$), also
$\mathfrak{x}pk=A+t-90^0$, $\mathfrak{y}pk=B+t-90^0$, $\mathfrak{z}pk=C+t-90^0$. Es wird demnach
\[
\begin{aligned}
x&=r\sin a\sin(A+t)\\
y&=r\sin b\sin(B+t)\\
z&=r\sin c\sin(C+t)
\end{aligned}
\]

Es bleibt uns jetzt noch übrig, die Grössen $a$, $A$ u.\,s.\,w., die nur von der
Lage der Bahn des Planeten, nicht von seinem jedesmaligen Orte in derselben
abhängig sind, aus der Neigung der Ebene dieser Bahn und der Länge des aufsteigenden
Knotens abzuleiten; wir bezeichnen jene mit $i$, diese mit $\AscendingNode$. Die
Betrachtung des Dreiecks $\mathfrak{x}pP$ gibt uns folgende drei Gleichungen:
\[
\begin{aligned}
\operatorname{cotang}\mathfrak{x}pP&=\frac{\sin pP\operatorname{cotang}\mathfrak{x}P-\cos pP\cos pP\mathfrak{x}}{\sin pP\mathfrak{x}}\\
\cos\mathfrak{x}p&=\cos pP\cos\mathfrak{x}P+\sin pP\sin\mathfrak{x}P\cos pP\mathfrak{x}\\
\sin\mathfrak{x}p&=\frac{\sin\mathfrak{x}P\sin pP\mathfrak{x}}{\sin\mathfrak{x}pP}
\end{aligned}
\]

Eben so geben die Dreiecke $\mathfrak{y}pP$, $\mathfrak{z}pP$ jedes drei ähnliche Gleichungen,
welche hier herzusetzen unnöthig ist, da man, um sie zu erhalten, in den drei
obigen nur $\mathfrak{x}$ mit $\mathfrak{y}$ und $\mathfrak{z}$ zu vertauschen hat. Nun ist
\[
pP=i,\quad pP\mathfrak{x}=90^0-\AscendingNode,\quad pP\mathfrak{y}=pP\mathfrak{z}=180^0-\AscendingNode
\]
Mit diesen und den übrigen Substitutionen werden unsre neun Gleichungen diese:
\[
\begin{aligned}
\operatorname{cotang}A&=-\cos i\operatorname{tang}\AscendingNode\\
\cos a&=\sin i\sin\AscendingNode\\
\sin a&=\frac{\cos\AscendingNode}{\sin A}
\end{aligned}
\]
\SourceLeaf{104}%
\[
\begin{aligned}
\operatorname{cotang}B&=\frac{-\sin i\operatorname{tang}\varepsilon+\cos i\cos\AscendingNode}{\sin\AscendingNode}\\
\cos b&=-\cos i\sin\varepsilon-\sin i\cos\varepsilon\cos\AscendingNode\\
\sin b&=\frac{\cos\varepsilon\sin\AscendingNode}{\sin B}\\
\operatorname{cotang}C&=\frac{\sin i\operatorname{cotang}\varepsilon+\cos i\cos\AscendingNode}{\sin\AscendingNode}\\
\cos c&=\cos i\cos\varepsilon-\sin i\sin\varepsilon\cos\AscendingNode\\
\sin c&=\frac{\sin\varepsilon\sin\AscendingNode}{\sin C}
\end{aligned}
\]

Die Unbestimmtheit, ob man $A$, $B$ und $C$ in den beiden ersten oder in
den beiden letzten Quadranten anzunehmen habe, wird man so entscheiden, dass
die Sinus von $a$, $b$ und $c$ positiv werden. Man nimmt also $A$ in den beiden
ersten Quadranten, wenn $\cos\AscendingNode$ positiv, $B$ und $C$ in eben denselben, wenn $\sin\AscendingNode$
positiv ist; in den entgegengesetzten Fällen aber in den beiden letzten Quadranten,
also $B$ und $C$ in den beiden ersten oder letzten, je nachdem $\AscendingNode$ in diesen
oder jenen liegt.

Die vierte, fünfte, siebente und achte dieser Gleichungen lassen sich durch
die Einführung von Hülfswinkeln noch bequemer einrichten. Dies kann auf
eine doppelte Weise geschehen:

\textit{Erstlich}, wenn man $\frac{\operatorname{tang}i}{\cos\AscendingNode}=\operatorname{tang}E$ und $\operatorname{tang}i\cos\AscendingNode=\operatorname{tang}F$ setzt, so wird
\[
\begin{aligned}
\operatorname{cotang}B&=\frac{\sin i\cos(E+\varepsilon)}{\sin\AscendingNode\cos\varepsilon\sin E}=\frac{\cos i\cos(E+\varepsilon)}{\operatorname{tang}\AscendingNode\cos\varepsilon\cos E}\\
\cos b&=-\frac{\cos i\sin(F+\varepsilon)}{\cos F}=-\frac{\sin i\cos\AscendingNode\sin(F+\varepsilon)}{\sin F}\\
\operatorname{cotang}C&=\frac{\sin i\sin(E+\varepsilon)}{\sin\AscendingNode\sin\varepsilon\sin E}=\frac{\cos i\sin(E+\varepsilon)}{\operatorname{tang}\AscendingNode\sin\varepsilon\cos E}\\
\cos c&=\frac{\cos i\cos(F+\varepsilon)}{\cos F}=\frac{\sin i\cos\AscendingNode\cos(F+\varepsilon)}{\sin F}
\end{aligned}
\]

\textit{Zweitens}, macht man $\frac{\operatorname{tang}\varepsilon}{\cos\AscendingNode}=\operatorname{tang}G$, und $\operatorname{tang}\varepsilon\cos\AscendingNode=\operatorname{tang}H$, so wird:
\[
\begin{aligned}
\operatorname{cotang}B&=\frac{\cos(G+i)}{\operatorname{tang}\AscendingNode\cos G}=\frac{\operatorname{tang}\varepsilon\cos(G+i)}{\sin\AscendingNode\sin G}\\
\cos b&=-\frac{\sin\varepsilon\sin(G+i)}{\sin G}=-\frac{\cos\AscendingNode\cos\varepsilon\sin(G+i)}{\cos G}\\
\operatorname{cotang}C&=\frac{\sin(H+i)}{\sin\AscendingNode\operatorname{tang}\varepsilon\cos H}=\frac{\sin(H+i)}{\operatorname{tang}\AscendingNode\sin H}\\
\cos c&=\frac{\cos\varepsilon\cos(H+i)}{\cos H}=\frac{\sin\varepsilon\cos\AscendingNode\cos(H+i)}{\sin H}
\end{aligned}
\]

\SourceLeaf{105}%
Es wird wohl der Mühe werth sein, noch einige Relationen zwischen
den Grössen $A$, $a$ u.\,s.\,w. zu entwickeln. Das sphärische Dreieck $\mathfrak{x}p\mathfrak{y}$
gibt $\cos\mathfrak{x}\mathfrak{y}=\cos\mathfrak{x}p\cos\mathfrak{y}p+\sin\mathfrak{x}p\sin\mathfrak{y}p\cos\mathfrak{x}p\mathfrak{y}$. Allein $\mathfrak{x}\mathfrak{y}=90^0$ und
$\mathfrak{x}p\mathfrak{y}=\mathfrak{x}pP-\mathfrak{y}pP=A-B$. Also
\[
\cos(A-B)=-\operatorname{cotang}a\operatorname{cotang}b
\]
Eben so geben die Dreiecke $\mathfrak{y}p\mathfrak{z}$, $\mathfrak{z}p\mathfrak{x}$
\[
\begin{aligned}
\cos(B-C)&=-\operatorname{cotang}b\operatorname{cotang}c\\
\cos(C-A)&=-\operatorname{cotang}c\operatorname{cotang}a
\end{aligned}
\]

Ferner wird in dem Dreiecke $\mathfrak{x}p\mathfrak{y}$, $\cos a=\cos p\mathfrak{y}\mathfrak{x}\sin b$, und in dem
Dreiecke $\mathfrak{y}p\mathfrak{z}$, $\sin\mathfrak{z}\mathfrak{y}p=\sin\mathfrak{y}p\mathfrak{z}\sin c$. Da nun $\mathfrak{z}\mathfrak{y}p+p\mathfrak{y}\mathfrak{x}=\mathfrak{z}\mathfrak{y}\mathfrak{x}=90^0$,
so hat man $\cos a=\sin b\sin c\sin\mathfrak{y}p\mathfrak{z}$, oder da $\mathfrak{y}p\mathfrak{z}=B-C$ ist
\[
\sin(B-C)=\frac{\cos a}{\sin b\sin c}
\]
Ganz auf ähnliche Art findet man
\[
\begin{aligned}
\sin(C-A)&=\frac{\cos b}{\sin c\sin a}\\
\sin(A-B)&=\frac{\cos c}{\sin a\sin b}
\end{aligned}
\]

Die Verbindung dieser Gleichungen mit den vorigen gibt noch
\[
\begin{aligned}
\operatorname{cotang}(A-B)&=-\frac{\cos a\cos b}{\cos c}=\operatorname{tang}(F+\varepsilon)\cos a\\
\operatorname{cotang}(B-C)&=-\frac{\cos b\cos c}{\cos a}\\
\operatorname{cotang}(C-A)&=-\frac{\cos c\cos a}{\cos b}=\operatorname{cotang}(F+\varepsilon)\cos a
\end{aligned}
\]
\[
\begin{aligned}
\cos a^2&=\operatorname{cotang}(A-B)\operatorname{cotang}(C-A)\\
\cos b^2&=\operatorname{cotang}(B-C)\operatorname{cotang}(A-B)\\
\cos c^2&=\operatorname{cotang}(C-A)\operatorname{cotang}(B-C)
\end{aligned}
\]
und auf ähnliche Art lassen sich die Quadrate der Sinus und Tangenten der Seiten
$a$, $b$, $c$ durch die Winkel $A-B$, $B-C$, $C-A$ darstellen.

\SourceLeaf{106}%
Um den Gebrauch dieser Formeln zu erläutern, wollen wir einige derselben
auf die Pallas anwenden, und dabei die neuesten [VII.] Elemente dieses Planeten
für 1803 zum Grunde legen. Wir setzen also
\[
\begin{aligned}
i&=34^0\;38'\;1''1\\
\AscendingNode&=172\;28\;13.7\\
\varepsilon&=23\;27\;55.8\quad\text{(mittlere Schiefe nach \textsc{Maskelyne} für 1803).}
\end{aligned}
\]

Mit diesen Elementen steht die Rechnung folgendermaassen (die den Logarithmen
beigesetzten n zeigen an, dass sie zu negativen Grössen gehören):

\begin{center}\begin{minipage}{\textwidth}\centering
\begin{tabular}{@{}lcl@{}}
$\log\cos i$ & $ .\quad .\quad .$ & $9.9152958$\\
$\log\operatorname{tang}\AscendingNode$ & $ .\quad .$ & $9.1211553\,\mathrm{n}$\\\hline
$\log\operatorname{cotang}A$ & $.$ & $9.0364511$
\end{tabular}\par\smallskip
 also $A=263^0\;47'\;35''4$
\end{minipage}\end{center}
\begin{center}\begin{minipage}{\textwidth}\centering
\begin{tabular}{@{}lcl@{}}
$\log\cos\AscendingNode$ & $ .\quad .\quad .$ & $9.9962390\,\mathrm{n}$\\
$\log\sin A$ & $ .\quad .\quad .$ & $9.9974467\,\mathrm{n}$\\\hline
$\log\sin a$ & $ .\quad .\quad .$ & $9.9987923$
\end{tabular}
\end{minipage}\end{center}
\begin{center}\begin{minipage}{\textwidth}\centering
\begin{tabular}{@{}lcl@{}}
$\log\sin i$ & $ .\quad .\quad .$ & $9.7545982$\\
$\log\sin\AscendingNode$ & $ .\quad .\quad .$ & $9.1173944$\\\hline
$\log\cos a$ & $ .\quad .\quad .$ & $8.8719926$
\end{tabular}\par\smallskip
 hieraus $a=85^0\;43'\;44''8$
\end{minipage}\end{center}
\begin{center}\begin{minipage}{\textwidth}\centering
\begin{tabular}{@{}lcl@{}}
$\log\operatorname{tang}i$ & $ .\quad .$ & $9.8393024$\\
$\log\cos\AscendingNode$ & $ .\quad .\quad .$ & $9.9962390\,\mathrm{n}$\quad also\\\hline
\end{tabular}
\[
\begin{gathered}
\log\operatorname{tang}E\quad .\quad .\quad 9.8430634\,\mathrm{n},\\
E=145^0\;8'\;2''4;\quad E+\varepsilon=168^0\;35'\;58''2\\
\log\operatorname{tang}F\quad .\quad .\quad 9.8355414\,\mathrm{n},\\
F=145\;35\;52.9;\quad F+\varepsilon=169\;3\;48.7
\end{gathered}
\]
\end{minipage}\end{center}
\begin{center}\begin{minipage}{\textwidth}\centering
\begin{tabular}{@{}lcl@{}}
$\log\cos i$ & $ .\quad .\quad .$ & $9.9152958$\\
Compl. $\log\operatorname{tang}\AscendingNode$ & & $0.8788447\,\mathrm{n}$\\
Compl. $\log\cos E$ & & $0.0859260\,\mathrm{n}$\\\hline
$\log\mathrm{const}.$ & $ .\quad .$ & $0.8800665$
\end{tabular}
\end{minipage}\end{center}
\begin{center}\begin{minipage}{\textwidth}\centering
\begin{tabular}{@{}lcl@{}}
$\log\cos(E+\varepsilon)$ & $.$ & $9.9913455\,\mathrm{n}$\\
Compl. $\log\cos\varepsilon$ & & $0.0374886$\\\hline
$\log\operatorname{cotang}B$ & $.$ & $0.9089006\,\mathrm{n}$
\end{tabular}\par\smallskip
 hieraus $B=172^0\;58'\;7''4$,
\end{minipage}\end{center}
\SourceLeaf{107}%
\begin{center}\begin{minipage}{\textwidth}\centering
\begin{tabular}{@{}lcl@{}}
$\log\mathrm{const}.$ & $ .\quad .$ & $0.8800665$\\
$\log\sin(E+\varepsilon)$ & $.$ & $9.2959318$\\
Compl. $\log\sin\varepsilon$ & & $0.3999023$\\\hline
$\log\operatorname{cotang}C$ & $.$ & $0.5759006$
\end{tabular}\par\smallskip
 $C=14^0\;52'\;12''5$
\end{minipage}\end{center}
\begin{center}\begin{minipage}{\textwidth}\centering
\begin{tabular}{@{}lcl@{}}
$\log\cos\varepsilon$ & $ .\quad .\quad .$ & $9.9625114$\\
$\log\sin\AscendingNode$ & $ .\quad .\quad .$ & $9.1173944$\\
Compl. $\log\sin B$ & & $0.9121791$\\\hline
$\log\sin b$ & $ .\quad .\quad .$ & $9.9920849$
\end{tabular}
\end{minipage}\end{center}
\begin{center}\begin{minipage}{\textwidth}\centering
\begin{tabular}{@{}lcl@{}}
$\log\sin\varepsilon$ & $ .\quad .\quad .$ & $9.6000977$\\
$\log\sin n$ & $ .\quad .\quad .$ & $9.1173944$\\
Compl. $\log\sin C$ & & $0.5906942$\\\hline
$\log\sin c$ & $ .\quad .\quad .$ & $9.3081863$
\end{tabular}
\end{minipage}\end{center}
\begin{center}\begin{minipage}{\textwidth}\centering
\begin{tabular}{@{}lcl@{}}
$\log\cos i$ & $ .\quad .\quad .$ & $9.9152958$\\
$\log\cos F$ & $ .\quad .$ & $9.9165035\,\mathrm{n}$\\\hline
 & & $9.9987923\,\mathrm{n}$
\end{tabular}
\end{minipage}\end{center}
\begin{center}\begin{minipage}{\textwidth}\centering
\begin{tabular}{@{}lcl@{}}
$\log\sin(F+\varepsilon)$ & $.$ & $9.2781142$\\
$\log\cos(F+\varepsilon)$ & $.$ & $9.9920399\,\mathrm{n}$\\\hline
$\log\cos b$ & $ .\quad .\quad .$ & $9.2769065$\\
$\log\cos c$ & $ .\quad .\quad .$ & $9.9908322$
\end{tabular}\par\smallskip
 also
$\begin{aligned}b&=79^0\;5'\;39''4\\c&=11\;43\;52.8\end{aligned}$
\end{minipage}\end{center}

Wenn man nur die Sinus von $a$, $b$, $c$ verlangt, so ist die Rechnung für
ihre Cosinus nicht nöthig, und man kann also auch den Hülfswinkel $F$ entbehren.
Will man aber auch $a$, $b$, $c$ selbst kennen, so dienen die Cosinus (wovon
nachher noch ein Gebrauch vorkommt) dazu, die Zweideutigkeiten, welche die
Sinus allein dabei übrig lassen, zu entscheiden. Auch geben sie dann, wenn
die Sinus näher bei 1 sind, eine schärfere Bestimmung, und zugleich eine Controle
für die Richtigkeit der Rechnung. Zu dieser letzten Absicht ist auch noch
der Umstand brauchbar, dass $\frac{\cos i}{\cos F}=\pm\sin a$ ist, wo das obere Zeichen gilt,
wenn $F$ mit $A$ zugleich in den beiden ersten oder letzten Quadranten liegt;
das untere, wenn $F$ in einer andern Hälfte des Umfanges angenommen ist als $A$.
\SourceLeaf{108}%
(Zur Entwickelung des Grundes davon dient die Bemerkung, dass $F$ im ersten
Falle mit dem Winkel $P\mathfrak{x}p$ einerlei, im zweiten $180^0$ davon verschieden
ist).

Die Grössen $\varepsilon$, $\AscendingNode$, $i$ sind Secularänderungen unterworfen: dasselbe wird
also auch der Fall mit den davon abhängigen $A$, $a$, $B$, $b$, $C$, $c$ sein. Sind die
jährlichen Änderungen von jenen bekannt, so können die Änderungen von $A$, $a$
u.\,s.\,w. durch leicht zu entwickelnde Differentialformeln berechnet werden, bei
welchen wir uns hier nicht aufhalten wollen. Man kann auch die Werthe von
$A$, $a$ u.\,s.\,w. für eine entferntere Epoche von neuem berechnen, und daraus ihre
jährlichen Änderungen ableiten.

Ausserdem erleiden diese Grössen wegen der Nutation noch periodische Änderungen,
die mit jedem Umlaufe der Mondsknoten wiederkehren. Da man nemlich
die geocentrische Lage des Planeten gegen den wahren Äquator verlangt, so
muss eigentlich für $\varepsilon$ nicht die mittlere, sondern die wahre Schiefe der Ecliptik,
und für $\AscendingNode$ die Entfernung des aufsteigenden Knotens vom wahren, nicht vom
mittlern Äquinoctialpunkte genommen werden. Die hieraus entspringenden periodischen
Änderungen können nach eben den Differentialformeln wie die Secularänderungen
berechnet, und in eine Tafel, deren Argument die Länge des Mondsknotens
ist, gebracht werden. Wenn man eine zahlreiche Menge geocentrischer
Örter für einen nicht zu grossen Zeitraum zu berechnen hat, wird man es in Ermangelung
einer solchen Tafel am bequemsten finden, für zwei Epochen zu Anfang
und Ende desselben die wahren Werthe von $A$, $a$ u.\,s.\,w. sogleich unmittelbar
aus den wahren Werthen von $\varepsilon$, $i$, $\AscendingNode$ zu berechnen, und für dazwischen liegende
Zeiten sie daraus durch einfache Interpolation abzuleiten. Ein Jahr hindurch
kann man ohne Bedenken diese Änderungen als gleichförmig ansehen.

Man könnte auch die von der Nutation abhängigen periodischen Änderungen
ganz übergehen, und sich der mittlern Werthe von $A$, $a$ u.\,s.\,w. bedienen:
dann müsste man aber auch bei der Erde für $\varepsilon$ die mittlere Schiefe der Ecliptik
gebrauchen, und von der Länge $\lambda$ die Nutation weglassen, um den Abstand
vom mittlern Äquinoctium zu haben. Der Erfolg davon ist sodann, dass man
die geocentrische Rectascension und Declination des Planeten in Beziehung auf
den \textit{mittlern} Äquator erhält, woraus man dann seine Lage gegen den \textit{wahren}
Äquator eben so ableitet, wie man den mittlern Ort eines Fixsterns durch Anbringung
der Nutation auf den scheinbaren reducirt.

\SourceLeaf{109}%
Wir haben jetzt nur noch einiges über die Perturbationen hinzuzufügen.
Die Störungen der Breite, von denen allein natürlich hier die Rede ist, sind
bei allen ältern Planeten so unbeträchtlich, dass man sie mit Recht ganz vernachlässigen
kann; blos bei der Ceres und Pallas wird es wegen der starken Neigung
der Bahnen dieser Planeten gegen die Jupitersbahn nothwendig, sie mit in
Rechnung zu nehmen. Es gibt dazu einen doppelten Weg. Man kann nemlich
entweder diejenigen Elemente, welche die Lage der Bahn bestimmen, die Neigung
und die Länge des Knotens, als veränderlich ansehen und ihre mittlern
Werthe durch periodische Gleichungen verbessern, oder auch gerade zu untersuchen,
wie viel der Planet aus der mittlern Ebene seiner Bahn herauszuweichen
durch fremde Kräfte genöthigt wird. Im ersten Falle wird man jene Änderungen
auch auf die Grössen $A$, $a$ u.\,s.\,w. übertragen, also diesen ausser den von
der Nutation abhängenden noch andere periodische Gleichungen beifügen, deren
Argumente mit denen für die Gleichungen der Neigung und der Länge des Knotens
übereinkommen werden. Dieses Verfahren ist jedoch bisher nicht üblich
gewesen. Bei der zweiten Methode hingegen werden die Störungstafeln die Perturbation
der heliocentrischen Breite angeben, welche aber eigentlich nichts anders
ist, als die heliocentrische Breite des Planeten über der mittlern Ebene seiner
Bahn. Es sei dieselbe $=\ObservedBeta$, gegen den Nordpol zu als positiv, gegen den
Südpol zu als negativ angesehen. In dem sphärischen Dreiecke $\mathfrak{x}pk$ ist also
die Seite $pk$ nicht wie vorhin $=90^0$, sondern $=90^0-\ObservedBeta$, folglich
\[
x=r\cos\mathfrak{x}k=r(\sin\ObservedBeta\cos a+\cos\ObservedBeta\sin a\sin(t+A))
\]
und eben so
\[
\begin{aligned}
y&=r(\sin\ObservedBeta\cos b+\cos\ObservedBeta\sin b\sin(t+B))\\
z&=r(\sin\ObservedBeta\cos c+\cos\ObservedBeta\sin c\sin(t+C))
\end{aligned}
\]

In so fern hier $\ObservedBeta$ höchstens nur einige Minuten betragen kann, wird man
$\cos\ObservedBeta=1$ und $\sin\ObservedBeta=\ObservedBeta$ setzen dürfen. Hieraus erhellt, dass man wegen der
Störungen zu den ohne sie gefundenen Werthen von $x$, $y$, $z$ nur noch die
Grössen $\ObservedBeta r\cos a$, $\ObservedBeta r\cos b$, $\ObservedBeta r\cos c$ hinzuzusetzen habe, wo $\ObservedBeta$ in Theilen des
Halbmessers ausgedrückt werden muss.

\medskip
[Handschriftliche Bemerkungen zu S.\,100.]
\[
\sin a^2=-\frac{\cos(B-C)}{\sin(A-B)\sin(C-A)}
\]
\SourceLeaf{110}%
\[
\operatorname{tang}a^2=-\frac{\cos(B-C)}{\cos(A-B)\cos(C-A)}
\]
\[
\cos a^2+\cos b^2+\cos c^2=1
\]
\[
\begin{aligned}
-\frac{\cos a}{\cos b}&=\operatorname{tang}(B-C)\cos=\frac{\operatorname{cotang}(C-A)}{\cos c}\\
&=\operatorname{tang}\AscendingNode'\ \text{in plano ipsius }Z
\end{aligned}
\]
\[
\begin{aligned}
\sin\AscendingNode'&=\sin b\sin(B-C)=\frac{\cos a}{\sin c}\\
\cos\AscendingNode'&=-\frac{\cos b}{\sin c}
\end{aligned}
\]

\begin{center}Differentialänderungen:\end{center}
\[
2\mathrm{d}a=-\frac{\mathrm{d}A-\mathrm{d}B}{\sin(A-B)\operatorname{cotang}b}-\frac{\mathrm{d}C-\mathrm{d}A}{\sin(C-A)\operatorname{cotang}c}
\]

Um umgekehrt aus $a$, $b$, $c$; $\AscendingNode$ und $i$ zu finden, dienen die Formeln
\[
\begin{aligned}
\cos a&=\sin i\sin\AscendingNode\\
-\cos b\cos\varepsilon-\cos c\sin\varepsilon&=\sin i\cos\AscendingNode\\
-\cos b\sin\varepsilon+\cos c\cos\varepsilon&=\cos i
\end{aligned}
\]
\[
\begin{aligned}
\sin a&=k\\
\cos b&=k\cos\theta\\
\cos c&=k\sin\theta
\end{aligned}
\]
\[
\begin{aligned}
\cos a&=\sin i\sin\AscendingNode\\
-k\cos(\theta-\varepsilon)&=\sin i\cos\AscendingNode\\
k\sin(\theta-\varepsilon)&=\cos i
\end{aligned}
\]
